\chapter{Cơ sở lý thuyết và nghiên cứu liên quan}

Chương này tổ chức văn liệu theo đối tượng được kiểm soát: tính khả thi của chuyến,
chất lượng của thông báo và lịch sử thay đổi thông báo. Phân biệt này xác định điều
RideBound kế thừa, điều cơ chế của khóa luận bổ sung và phạm vi cần kiểm tra bằng thực nghiệm.

\section{Điều phối động và chất lượng phục vụ}

\subsection{Bài toán và các ràng buộc cơ bản}

Bài toán Dial-a-Ride Problem (DARP) ghép việc gán hành khách cho xe với việc sắp xếp
các điểm đón và trả \cite{cordeau2007}. Trong dạng động, yêu cầu chưa biết trước toàn bộ
mà xuất hiện trong lúc xe đang vận hành. Một yêu cầu gồm điểm đón, điểm trả, thời điểm
yêu cầu và số chỗ cần. Trạng thái hiện tại gồm vị trí xe, khách đã nhận, khách đang trên
xe và các điểm dừng đã cam kết.

Tính khả thi vật lý thường gồm sức chứa, đón trước trả, cửa sổ đón và giới hạn thời gian
trên xe. Giữ nguyên một tiền tố tuyến còn tránh đổi những chặng đã thực hiện hoặc đang
không thể đổi. Chúng trả lời câu hỏi ``xe có thể thực hiện kế hoạch này hay không''.
Một ràng buộc lời hứa trả lời thêm câu hỏi ``kế hoạch này có được phép thay đổi thông tin
đã báo cho khách hay không''. Hai câu hỏi cần được kiểm tra riêng vì có thể thất bại vì
những lý do khác nhau.

Chẳng hạn, hai phương án đón khách ở phút thứ ba và phút thứ sáu đều thỏa thời gian chờ
tối đa mười phút. Nếu khách đã nhận thông báo đón ở phút thứ ba, phương án thứ hai thay
đổi thông báo dù vẫn khả thi vật lý. Ngược lại, giữ nguyên tuyến trong một đợt tăng thời
gian đi có thể làm xe đón sau phút thứ mười. Một cổng cam kết chấp nhận việc giữ tuyến
không khôi phục được cửa sổ đón đã vi phạm ấy.

\subsection{Sinh phương án và chọn phương án}

Alonso-Mora và cộng sự \cite{alonsomora2017} sử dụng đồ thị
Request--Trip--Vehicle để biểu diễn các nhóm yêu cầu và xe có thể phục vụ nhóm đó,
rồi giải bài toán gán. Trong thực nghiệm taxi New York của họ, 3.000 xe sức chứa bốn
người phục vụ 98\% nhu cầu, với thời gian chờ trung bình 2,8 phút và độ trễ chuyến trung
bình 3,5 phút. Đây là kết quả trong mô hình và điều kiện của bài; khóa luận không dùng
98\% như một mốc tỉ lệ phục vụ mà mạng dẫn xuất nhỏ phải đạt.

Một dạng bài toán chọn phương án có thể viết như sau. Gọi $\Gamma$ là tập phương án
xe--chuyến đã sinh, $x_T$ là biến chọn phương án $T$ và $z_r$ là biến từ chối yêu cầu $r$:

\begin{equation}
\min_{x,z}\quad \sum_{T\in\Gamma}\mathrm{cost}(T)x_T
  +p_{\mathrm{KO}}\sum_{r\in R}z_r
\end{equation}
\begin{align}
z_r+\sum_{T:\,r\in\mathrm{req}(T)}x_T&=1 &&\forall r\in R,\\
\sum_{T:\,\mathrm{veh}(T)=v}x_T&\leq1 &&\forall v\in V,\\
x_T,z_r&\in\{0,1\}.&&
\end{align}

Mô hình này buộc mỗi yêu cầu hoặc được phục vụ bởi một phương án, hoặc chịu hình phạt
từ chối; mỗi xe được chọn tối đa một phương án. Nó minh họa sự tách biệt giữa sinh và chọn.
Chương 3 mô tả bộ sinh và thứ tự mục tiêu thực tế của RideBound, gồm ưu tiên số khách
nhận rồi mới so chi phí. Vì vậy, công thức trên là kiến thức nền, không phải lời khẳng
định rằng mọi hệ thống dùng cùng một hàm mục tiêu.

Chi phí có thể xét cả khách mới và phần thiệt hại tăng thêm của khách cũ. Fielbaum,
Tirachini và Alonso-Mora \cite{fielbaum2022scale} phân tích các chi phí của người dùng
và nhà điều hành trong mô hình ghép chuyến. Việc đưa thiệt hại của khách cũ vào hàm mục
tiêu là một tiền lệ quan trọng. Tuy nhiên, ưu tiên mềm được tổng hợp giữa nhiều khách
khác với ngân sách cứng kiểm riêng từng khách. Một phương án có chi phí tổng thấp vẫn
có thể gây thay đổi lớn cho một người; ngân sách theo hành khách kiểm tra đúng trường
hợp đó trước khi chọn.

\section{Đo thay đổi và giới hạn lời hứa}


\subsection{Chênh lệch kết quả và lịch sử cập nhật}

Fielbaum và Alonso-Mora \cite{fielbaum2020} phân biệt bất ổn trong một chuyến với
bất ổn giữa các lần thực hiện cùng yêu cầu. Với một chuyến, họ so thời gian thực tế với
thông báo đầu và đếm số lần dự báo tăng. Phép so hai đầu về giờ chờ có dạng
$\Delta t_w=t_w-t_w^1$; các phần âm được cắt về không khi đo sự xấu đi.

Một thước đo hai đầu phù hợp khi câu hỏi là độ lệch cuối cùng so với thông báo đầu.
Nó không mô tả đầy đủ độ dài chuỗi cập nhật: một dự báo tăng rồi giảm về mốc ban đầu
có chênh lệch cuối bằng không nhưng tổng biến thiên dương. Đếm số lần cập nhật cung
cấp thông tin bổ sung, song hai chuỗi cùng số lần có thể có độ lớn rất khác nhau.
Khóa luận vì thế báo cả chỉ số kết quả và chỉ số lịch sử, thay vì xem một chỉ số thay
thế toàn bộ các chỉ số còn lại.

Đại lượng thực tế chỉ xác định sau chuyến; dự báo của một phương án có thể xác định
trước khi chọn phương án. Sự khác biệt này đòi hỏi xác định rõ đại lượng nào được
kiểm ở mỗi thời điểm. Nó không có nghĩa là hiện tượng trong thước đo kết quả không
thể kiểm soát. Chính phần 6 của bài năm 2020 đã siết thời gian chờ tối đa xuống dự báo
và cấm từ chối lại khách đã nhận, trong kịch bản kiểm soát hoàn toàn. Tỉ lệ từ chối ở
bảng 5 tăng từ 34,5\% lên 45,6\%, tức 11,1 điểm phần trăm, trong hệ của các tác giả.

Phần mở rộng được xét trong RideBound là dùng mốc chiếu giữ tuyến tại cùng trạng thái
để phân biệt cập nhật gắn với lựa chọn kế hoạch và cập nhật khi kế hoạch không đổi.
Bài năm 2020 loại tắc đường và tai nạn khỏi phạm vi phân tích của nó. Vì vậy, đối
chiếu với RideBound dưới trôi dạt cần giữ khác biệt mô hình này; không thể lấy kết
quả của hai hệ như một cuộc so sánh trực tiếp về hiệu quả.

\subsection{Bảo vệ theo độ chùng và theo thông báo đầu}

Elmi và cộng sự \cite{elmi2026} dùng chỉ số shareability để dự báo khả năng hành khách
bị đổi thời gian, rồi phân bổ độ chùng có chọn lọc: nhóm chỉ số thấp không được tăng
thời gian, nhóm chỉ số cao giữ độ chùng. Đối chứng của họ điều chỉnh cùng một tỉ lệ độ
chùng cho mọi khách. Cả hai đều là cơ chế kiểm soát theo từng khách, không chỉ một
hình phạt trung bình trên toàn đội xe.

Điều này thu hẹp ranh giới của khóa luận. Một trần áp cho mọi khách không tự nó tạo
tính mới. Khác biệt cần kiểm tra là ngân sách đang chặn độ lệch từ lúc gán, tổng độ
lớn các cập nhật, hay phần do quyết định sau khi đối chiếu với phương án giữ tuyến.
Trong mạng có thời gian đi cố định và một bộ sinh chỉ chèn thêm, các đại lượng này
có thể gần nhau. Khi thời gian đi trôi dạt hoặc dự báo đổi chiều, chúng có thể tách
nhau. Chương 3 phát biểu điều kiện; chương 4 kiểm tra mức tách quan sát được.

Schulz và Pfeiffer \cite{schulz2026} giới hạn thời điểm đến điểm trả của khách đã
chấp nhận bằng thông báo lúc nhận cộng một khoảng chung. Một dạng ký hiệu là
\begin{equation}
p_t-p_0\leq\kappa,
\end{equation}
với $p_0$ là giờ trả đã công bố khi nhận khách. Đây là tiền lệ trực tiếp cho cổng
hạn chót một chiều. Nó giới hạn mức muộn hiện tại, không cộng độ lớn mọi lần sửa
thông báo và không chặn chiều sớm. Khóa luận cài đặt một đối chứng cùng dạng ràng
buộc này, không tái lập toàn bộ thuật toán, mô hình chấp nhận hoặc dữ liệu của bài.

Haferkamp \cite{haferkamp2024} nghiên cứu nhiều lựa chọn giờ đón và sự đánh đổi giữa
khả năng khách nhận đề nghị với công sức định tuyến. Sau khi khách chọn đề nghị, giờ đón
và phương án gán đã thống nhất trở thành ràng buộc. Các lần sửa tuyến phải giữ giờ đến
trong cửa sổ phục vụ tương ứng, có cho phép một độ trễ đón nhỏ.
Sivagnanam và cộng sự
\cite{sivagnanam2026} kết hợp xác nhận nhanh yêu cầu đặt trước với tối ưu tiếp tuyến
cho khách đã nhận. Schmidt và cộng sự \cite{schmidt2026} xét cam kết giờ đến trong
một bài toán người bán hàng có nhu cầu phụ thuộc lời chào. Những cơ chế này cho thấy
công bố thông tin, giữ khách đã nhận và kiểm soát giờ phục vụ đều đã có tiền lệ.
Chúng không được dùng như đối chứng số trong thực nghiệm RideBound vì khác bài toán
và khác mô hình nhu cầu.

\subsection{Nhất quán, nhiều kế hoạch và phương án tham chiếu}

Tính nhất quán giờ phục vụ đã được xem xét trong các bài toán định tuyến nhiều kỳ
\cite{kovacs2014}. Tellez và cộng sự \cite{tellez2022} xét mức nhất quán thời gian
cho từng hành khách qua các ngày và sự đánh đổi với chi phí. Đối tượng đó khác lịch
sử sửa một thông báo trong cùng chuyến. Khác biệt về đối tượng cho phép chọn phép
đo khác; nó không cho phép gọi tính nhất quán thời gian nói chung là một phát hiện mới.

Ackermann và Rieck \cite{ackermann2025} giữ nhiều kế hoạch để hỗ trợ chèn yêu cầu động
và bàn cách tránh cam kết sớm vào một kế hoạch. Vì vậy, nhiều kế hoạch và trì hoãn cam
kết không thuộc phần tính mới của RideBound. Tương tự, Wang và Fielbaum
\cite{wang2026walking} dùng lịch tham chiếu và cơ chế quay lại trong việc chọn điểm
trả có đi bộ. Bảo vệ không làm xấu một cập nhật đã có tiền lệ; đối tượng cần xét ở
đây là cách ghi ngân sách xuyên suốt lịch sử và xử lý vi phạm dưới trôi dạt.

Bảng~\ref{bang:ranhgioivanlieu} tóm tắt đối tượng được kiểm soát và ranh giới so sánh.

\begin{table}[htbp]
\centering
\caption{Đối tượng kiểm soát của các cơ chế gần nhất}
\label{bang:ranhgioivanlieu}
\small
\begin{tabular}{@{}>{\raggedright\arraybackslash}p{3.2cm}>{\raggedright\arraybackslash}p{5.4cm}>{\raggedright\arraybackslash}p{5.4cm}@{}}
\toprule
\textbf{Cơ chế} & \textbf{Đối tượng chính} & \textbf{Phân biệt cần giữ}\\
\midrule
Fielbaum và Alonso-Mora & Thời gian so với thông báo đầu; siết dự báo giờ chờ & Đã có kiểm
soát trước quyết định; không nghiên cứu trôi dạt giao thông\\[3pt]
Elmi và cộng sự & Độ chùng từng khách, phân bổ theo shareability & Bảo vệ theo khách đã có;
không đồng nhất độ chùng với sổ cái hành vi\\[3pt]
Schulz--Pfeiffer & Hạn chót điểm trả từ lúc nhận khách & Trần muộn một chiều khác tổng
biến thiên qua nhiều lần sửa\\[3pt]
RideBound & Sổ cái lịch sử theo khách, có mốc chiếu giữ tuyến & Giới hạn phần do quyết định
khác bảo đảm toàn bộ kết quả thực tế\\
\bottomrule
\end{tabular}
\end{table}

\section{Ổn định lịch trình và cách chọn thước đo}

Dao động giữa các lần lập kế hoạch không chỉ xuất hiện trong điều phối xe. Văn liệu
hoạch định sản xuất đã nghiên cứu nervousness và việc đóng băng một phần lịch.
Sridharan, Berry và Udayabhanu \cite{sridharan1987} xem xét các quyết định đóng băng
lịch sản xuất; bài năm 1988 \cite{sridharan1988} nghiên cứu cách đo độ ổn định dưới
chân trời cuốn chiếu. Hai bài có vai trò khác nhau và cần được trích riêng.

Khóa luận mượn tư tưởng đo thay đổi qua nhiều lần lập kế hoạch và kiểm tra độ nhạy với
trọng số theo chân trời. Nó không giả định bài toán sản xuất và bài toán vận tải có
cùng chi phí, cùng tính khả thi hoặc cùng đường cong đánh đổi. Việc dùng thước đo
được điều chỉnh cho lời hứa của hành khách là một lựa chọn đo lường, chưa chứng minh
sự tương đương giữa hai lĩnh vực.

Đối với một chuỗi dự báo $p_0,\ldots,p_T$, ba đại lượng cơ bản là chênh lệch cuối
$p_T-p_0$, số lần thay đổi và tổng biến thiên $\sum_{t=1}^T|p_t-p_{t-1}|$.
Trọng số thời gian hoặc phép bình phương nhấn mạnh những khía cạnh khác nhau của
chuỗi. Chương 4 dùng các biến thể này như kiểm tra độ nhạy của kết luận. Nếu thứ hạng
không đổi thì kết luận ổn định hơn đối với các phép đo đã thử; điều đó vẫn không
chứng minh mọi phép đo hoặc mọi đánh giá của người dùng sẽ cho cùng thứ hạng.

\section{Độ tin cậy và giá trị đối với hành khách}

Nghiên cứu lựa chọn dịch vụ cung cấp động lực cho việc đo thông báo, nhưng không trực
tiếp đặt được ngân sách số cho RideBound. Alonso-González và cộng sự
\cite{alonsogonzalez2020} nghiên cứu đánh đổi thời gian, độ tin cậy và chi phí bằng
thực nghiệm lựa chọn công bố tại Hà Lan. Trong mô hình của họ, giá trị độ tin cậy
tương đối được xác định qua độ biến thiên của thời gian; nó không phải đơn giá cho
mỗi giây sửa ETA.

Geržinič và cộng sự \cite{gerzinic2023} xét chênh lệch giữa thời gian chờ dự kiến và
thời gian chờ trải nghiệm trong các lựa chọn dịch vụ ride-hailing lặp lại. Cảm nhận
thời gian chờ ngoài dự kiến và kinh nghiệm trước ảnh hưởng tới lựa chọn. Đối tượng
đo này cũng khác tổng độ lớn cập nhật dự báo giờ trả của một chuỗi quyết định.
Không có phép chuyển đổi được kiểm chứng từ một trong hai mô hình đó sang phúc lợi
của những hành khách trong mô phỏng RideBound.

Bansal và cộng sự \cite{bansal2019} nghiên cứu cách hiển thị thời gian chờ dự báo và
mô phỏng chiến lược công bố. Kết quả về tăng mức nhận chuyến trong bài là kết quả
của mô hình lựa chọn và tối ưu thông tin, không phải bằng chứng rằng một sổ cái cam
kết bất kỳ sẽ làm tăng nhu cầu. Khóa luận vì thế giữ tỉ lệ phục vụ và biến động lời
hứa như hai đầu ra riêng, chưa tổng hợp chúng thành một chỉ số tiền tệ hoặc hài lòng.

\section{Cơ sở ghi sổ và phạm vi so sánh}

Các nhóm văn liệu trên chuẩn bị hai lựa chọn sẽ được định nghĩa trong chương 3.
Lựa chọn thứ nhất ghi toàn bộ biến động theo cơ sở thông báo nhìn thấy. Lựa chọn thứ hai
dùng mốc chiếu giữ tuyến tại cùng trạng thái để chỉ ghi phần chênh lệch của phương án được chọn.
Cả hai có thể dùng lịch sử tích lũy, nhưng chúng cam kết với khách về những đại lượng
khác nhau. Điều này giải thích vì sao một chính sách có ngân sách được tôn trọng vẫn
có thể không giữ được một trần trên biến động nhìn thấy.

Khóa luận dùng các nguồn gần nhất để xác định cơ chế đối chứng và ranh giới của câu
hỏi, không để khẳng định chưa từng có mọi thành phần tương tự. Những so sánh định
lượng trong chương 4 đều được thực hiện trên cùng bộ thực thi và đầu vào. Chúng đánh giá
các dạng cổng và cấu hình đã hiện thực; chúng chưa xác định thứ hạng giữa RideBound
với toàn bộ thuật toán của các bài gốc. Phần giá trị cần kiểm tra còn lại là tính
nhất quán của sổ cái, hệ quả của từng cơ sở ghi sổ và sự đánh đổi dịch vụ trong phạm
vi thực nghiệm đã khóa.
